\chapter{Register-level interface \& internal state encodings}
\label{ch:regs}

\begin{fnwarn}[Real SPI opcode map exists; state encodings below are
per-module reference]
Ch.~\ref{ch:host} now documents V2's real, physical SPI opcode map
(\code{WRITE\_JOB}/\code{WRITE\_MEM}/\code{READ\_MEM}/\code{STATUS}/
\code{RESET}) --- this chapter's own node-registration field layout
below remains the logical field reference (repeated here for quick
reference). The \textbf{internal FSM state encodings} below are useful
for simulation-level debug; \S\S\ref{ch:regs}'s Dependency
Manager/Neural Director tables are shared by every V2 architecture
(unchanged between the PSRAM-era and current SDRAM boards). The Memory
Manager and Neural Processor tables were captured from the PSRAM-era
\code{memory\_manager.v}/\code{neural\_processor.v} pairing (ch.~\ref{ch:arch})
--- the current SDRAM board's \code{nms\_memory\_manager\_stream\_wide.v}
implements the same functional handshake (prefetch $\to$ stream $\to$
write-back $\to$ done) against the SDRAM backend instead of PSRAM, but
its own internal state encoding was not re-transcribed into this table.
\end{fnwarn}

\section{Node registration fields (quick reference)}
See ch.~\ref{ch:host} for the full field-level description.
\code{reg\_node\_id}, \code{reg\_required}, \code{reg\_producer\_ids},
\code{reg\_x\_base}, \code{reg\_w\_base}, \code{reg\_n\_tiles},
\code{reg\_result\_addr} --- valid/ready handshake, \code{reg\_ready}
gated on the target node id's table slot being \code{EMPTY}.

\section{Dependency Manager node state (\texttt{node\_state})}
\begin{tabularx}{\textwidth}{C{1.4cm} L{2.6cm} Y}
\toprule
\rowh \thd{Value} & \thd{Name} & \thd{Meaning} \\
\midrule
\op{2'd0} & \code{ST\_EMPTY} & Table slot free; \code{reg\_ready} asserted for this node id. \\
\rowa \op{2'd1} & \code{ST\_WAITING} & Registered, at least one producer not yet resolved. \\
\op{2'd2} & \code{ST\_READY} & All producers resolved; eligible for dispatch. \\
\rowa \op{2'd3} & \code{ST\_DISPATCHED} & Handed to the Director; \textbf{terminal} (\S\ref{ch:sched}). \\
\bottomrule
\end{tabularx}

\section{Neural Director state (\texttt{dir\_state})}
\begin{tabularx}{\textwidth}{C{1.4cm} L{2.6cm} Y}
\toprule
\rowh \thd{Value} & \thd{Name} & \thd{Meaning} \\
\midrule
\op{4'd0} & \code{DIR\_IDLE} & Reset/startup. \\
\rowa \op{4'd1} & \code{DIR\_SCAN\_READY} & Checking whether a queued job and a free slot both exist. \\
\op{4'd2} & \code{DIR\_ALLOCATE} & Dispatching the head-of-queue job to the first free slot. \\
\rowa \op{4'd3} & \code{DIR\_ERROR} & Recoverable only via reset (an isolated fault never blocks other slots). \\
\bottomrule
\end{tabularx}

\section{Memory Manager state (\texttt{state})}
\begin{tabularx}{\textwidth}{C{1.4cm} L{3.0cm} Y}
\toprule
\rowh \thd{Value} & \thd{Name} & \thd{Meaning} \\
\midrule
\op{3'd0} & \code{MM\_IDLE} & Waiting for \code{job\_start}. \\
\rowa \op{3'd1} & \code{MM\_PREFETCH\_FIRST} & Waiting for tile~0's activation \emph{and} weight halves to both arrive. \\
\op{3'd2} & \code{MM\_STREAM} & Presenting tiles to the Neural Processor, double-buffering the next one. \\
\rowa \op{3'd3} & \code{MM\_WAIT\_RESULT} & Last tile handed off; waiting for the Neural Processor's own result. \\
\op{3'd4} & \code{MM\_WRITE\_RESULT} & Issuing the real PSRAM word write for the INT8 result. \\
\rowa \op{3'd5} & \code{MM\_DONE} & Waiting for the write's own \code{mem\_ready}; then pulses \code{job\_done}. \\
\bottomrule
\end{tabularx}

\section{Neural Processor state (\texttt{np\_state})}
\begin{tabularx}{\textwidth}{C{1.4cm} L{2.8cm} Y}
\toprule
\rowh \thd{Value} & \thd{Name} & \thd{Meaning} \\
\midrule
\op{4'd0} & \code{NP\_IDLE} & No job in flight. \\
\rowa \op{4'd1} & \code{NP\_LOAD\_JOB} & Latching \code{job\_bias}/\code{job\_activation}, clearing the accumulator. \\
\op{4'd2} & \code{NP\_WAIT\_OPERANDS} & Consuming tiles as they arrive (absorbs the per-tile MAC/accumulate/next-tile sequence). \\
\rowa \op{4'd3} & \code{NP\_FINISH} & Draining the pipeline after \code{tile\_last}. \\
\op{4'd4} & \code{NP\_WRITE\_RESULT} & Result available for the Memory Manager to consume. \\
\rowa \op{4'd5} & \code{NP\_DONE} & Job complete. \\
\op{4'd6} & \code{NP\_ERROR} & Reachable only via an unreachable \code{default} case --- isolated per-processor, never blocks other slots. \\
\bottomrule
\end{tabularx}

\section{Slot Memory Arbiter owner encoding}
\code{owner} is \code{0} for ``no port granted'', or (port index $+1$)
for the currently-granted port --- indices \code{0..N\_SLOTS-1} are the
per-slot Memory Managers' own weight/write-back traffic; index
\code{N\_SLOTS} is the shared Activation Cache's own traffic.
